
\element{general}{
    \begin{question}{classe}
        Quels sont les objets fonctionnels présents dans le système de palettisation automatique ? Pour chacun, indiquez :
        \begin{itemize}
            \item Un court descriptif de son rôle dans le système
            \item Les missions qu'il peut recevoir
            \item Les états possibles
            \item Son interface électrique
            \item Son interface logicielle (missions qu'il peut donner)
        \end{itemize}
        Votre réponse devra s'appuyer sur un schéma montrant les échanges entre les objets du système. 
        \evaluationProfLignes{10}{12}
    \end{question}
}

\element{partie1}{
    \begin{question}{station}
        Proposer un nom pour chaque variable globale associées aux différents capteurs et actionneurs de la station d'alimentation. \\
        \evaluationProfLignes{2}{4}
    \end{question}
}

\element{partie1}{
    \begin{question}{appel}
        Donner la ligne permettant d'appeler un bloc fonctionnel \lstinline{oStation} de type \lstinline{CStationAlimentation} dans le programme principal. \textit{(On suppose que sa déclaration a bien été faite dans le contexte du programme.)}\\
        \evaluationProfLignes{2}{3}
    \end{question}
}

\element{partie1}{
    \begin{question}{etats}
        Proposez une structure `STRUCT` décrivant les états possibles de la station d'alimentation. \textit{Votre réponse sera la déclaration de la structure en langage ST.}
        \evaluationProfLignes{5}{5}
    \end{question}
}

\element{partie1}{
    \begin{question}{missions}
        Proposez une structure décrivant les missions qu'elle peut recevoir. \textit{Votre réponse sera la déclaration de la structure en langage ST.}
        \evaluationProfLignes{5}{5}
    \end{question}
}

\element{partie1}{
    \begin{question}{missionsDeclaration}
        Donner le contexte (déclaration des variables) de la classe \lstinline{CStationAlimentation}.\\
        \textit{Votre réponse contiendra la déclaration de la classe, les variables d'entrées, de sorties ainsi que les variables représentant les états et les missions.}\\
        \evaluationProfLignes{2}{5}
    \end{question}
}

\element{partie1}{
    \begin{question}{getSet}
        Pour chacun des membres décrits (états et missions), définissez les propriétés `GET` et `SET` permettant de les manipuler. \\
        \textit{Votre réponse contiendra le corps (programme) de ces accesseurs.}\\
        \textit{Vous indiquerez également celles qui vous semblent inutiles en justifiant votre réponse.}
        \evaluationProfLignes{5}{5}
    \end{question}
}

\element{partie1}{
    \\
    Pour la station d'alimentation, on propose de coder sa mission principale (concernant le chargement de la zone tampon) à l'aide d'une action codée en SFC. 

    Le nombre de rangées à charger est stockée dans une variable d'instance \lstinline{m_nbRangees} de type \lstinline{UINT}.

    \begin{question}{grafcet_station}
        Proposer un programme SFC décrivant la mission principale de la station d'alimentation.\\
        \textbf{Ce programme ne gérera pas la mise en quinconce des rangées.}\\
        \evaluationProfLignes{5}{15}
    \end{question}
}

\element{partie1}{
    \begin{question}{setoutput}
        Proposer le corps de la fonction permettant de mettre à jour les sorties de la station d'alimentation.
        
        \textit{Bonus : Il est possible de gérer la mise en quinconce des rangées directement dans cette partie. Proposer la ligne de commande du vérin de butée mobile permettant de le faire.}\\
        \evaluationProfLignes{5}{6}
    \end{question}
}



